% Βοηθητικές εντολές του κεφαλαίου Γ2
% \binhist{ύψη}{άκρα κλάσεων}{μέγιστο ύψος}{τίτλος άξονα x}: ιστόγραμμα με κλάσεις ίσου πλάτους
\providecommand{\binhist}[4]{\begin{tikzpicture}[x=0.95cm,y=0.42cm]
  \foreach \y in {0,...,#3}{\draw[gray!40] (0,\y)--(5,\y);}
  \foreach \h [count=\i from 0] in {#1}{\draw[fill=gblue!35,draw=gblue,thick] (\i,0) rectangle ({\i+1},\h);}
  \foreach \l [count=\i from 0] in {#2}{\node[below,font=\scriptsize] at (\i,0) {\l};}
  \foreach \y in {0,...,#3}{\node[left,font=\scriptsize] at (0,\y) {\y};}
  \node[below,font=\small] at (2.85,-0.75) {#4};
  \node[rotate=90,font=\small] at (-0.9,{#3/2}) {συχνότητα};
\end{tikzpicture}}
\kefalaio{Μέτρα θέσης – Μεταβλητότητα}{Δύο παρέες μπορούν να έχουν τον ίδιο μέσο όρο και εντελώς διαφορετική ιστορία}{%
\begin{tabular}{@{}l@{\hspace{1em}}p{0.82\linewidth}@{}}
2.1 & Μέση τιμή, διάμεσος\\[2pt]
2.2 & Εύρος, κορυφές, απόμακρες τιμές, μεταβλητότητα των δεδομένων\\[2pt]
 & Επανάληψη του κεφαλαίου\\
\end{tabular}}
\paragrafos{Μέση τιμή και διάμεσος}{Γ.2.1, «Μέση τιμή, διάμεσος»}

\begin{anagke}
Η Ρήνη περιμένει κάθε πρωί το λεωφορείο για το σχολείο και μια εβδομάδα κρατά σημειώσεις: έχασε $5$, $6$, $4$, $5$, $7$, $5$ και $33$ λεπτά. Την Παρασκευή υπήρχε απεργία και το λεωφορείο καθυστέρησε πολύ. «Συνήθως περιμένω γύρω στα πέντε λεπτά», λέει η Ρήνη. Ο Ηλίας προσθέτει τους αριθμούς και τους μοιράζει στις μέρες: «Εγώ βρίσκω περίπου $9{,}3$ λεπτά». Ποιος έχει δίκιο;

Η κυρία Αλεξάνδρα ρωτά: «Και οι δύο έχετε δίκιο; Τι ακριβώς ψάχνει ο καθένας σας;». Επτά αριθμοί είναι δύσκολο να τους θυμάται κανείς. Χρειαζόμαστε \emph{έναν} αριθμό που να τους αντιπροσωπεύει και να μας λέει «πού βρίσκονται» τα δεδομένα. Αυτό όμως μπορεί να γίνει με δύο τρόπους.

Σε αυτή την παράγραφο μαθαίνουμε δύο μέτρα θέσης, τη μέση τιμή και τη διάμεσο, τα υπολογίζουμε από λίστα, από πίνακα συχνοτήτων και από ομαδοποιημένα δεδομένα, και καταλαβαίνουμε πότε προτιμούμε το ένα και πότε το άλλο.
\end{anagke}

\begin{lexi}{μέτρο θέσης, μέση τιμή, διάμεσος}{ονόματα και σημασίες}
Το \lexiword{μέτρο θέσης} δείχνει «πού βρίσκεται» το κέντρο των δεδομένων. Η \lexiword{μέση τιμή} λέγεται και \emph{μέσος όρος}: είναι ο αριθμός που θα είχαμε αν μοιράζαμε τα πάντα εξίσου. Η \lexiword{διάμεσος} είναι η τιμή που κάθεται «στη μέση» (δια-μέσος) όταν βάλουμε τα δεδομένα στη σειρά.
\end{lexi}

\begin{metafrasi}
Προσθέτω όλες τις τιμές και τις μοιράζω εξίσου. & μέση τιμή\\
Βάζω τις τιμές στη σειρά και παίρνω αυτή που είναι στη μέση. & διάμεσος\\
Μετρήσεις που κάπου τις έχω γράψει πολλές φορές. & πίνακας συχνοτήτων\\
Μετρήσεις που τις έχω χωρίσει σε κλάσεις (διαστήματα). & ομαδοποιημένα δεδομένα\\
\end{metafrasi}

\tmima{Η θεωρία}

\begin{kanonas}{μέση τιμή}
$$\text{μέση τιμή}=\frac{\text{άθροισμα των δεδομένων}}{\text{πλήθος των δεδομένων}}$$
\end{kanonas}

Η μέση τιμή είναι σαν σημείο ισορροπίας. Αν τοποθετήσεις τα δεδομένα $8, 9, 12, 15, 21$ πάνω σε μια αριθμογραμμή, σαν βαράκια σε μια ζυγαριά, το σημείο ισορροπίας βρίσκεται στο $13$: ο άξονας δεν γέρνει ούτε δεξιά ούτε αριστερά.

\begin{center}
\begin{tikzpicture}[x=0.5cm]
\draw[-{Stealth},thick] (6,0)--(23,0);
\foreach \x in {6,8,...,22}{\draw (\x,0.1)--(\x,-0.1) node[below,font=\scriptsize]{$\x$};}
\foreach \x in {8,9,12,15,21}{\fill[gblue] (\x,0.3) circle (3.5pt);}
\fill[gwarmD] (13,-0.55)--(12.65,-1.15)--(13.35,-1.15)--cycle;
\node[font=\small,gwarmD] at (13,-1.6) {μέση τιμή $=13$};
\end{tikzpicture}
\end{center}

\begin{idiotita}{τι ξέρουμε για τη μέση τιμή}
\begin{itemize}[leftmargin=1.2em,itemsep=0.3ex,topsep=2pt]
\item Είναι πάντα ανάμεσα στη μικρότερη και τη μεγαλύτερη τιμή.
\item Δεν είναι απαραίτητο να είναι μία από τις τιμές των δεδομένων (και μπορεί να είναι δεκαδικός).
\item Αν $\nu$ δεδομένα έχουν μέση τιμή $m$, τότε το άθροισμά τους είναι $\nu\cdot m$.
\item Την επηρεάζουν πολύ οι πολύ μεγάλες ή οι πολύ μικρές τιμές.
\end{itemize}
\end{idiotita}

\begin{kanonas}{διάμεσος}
Βάζω τα δεδομένα σε \textbf{αύξουσα σειρά}. Αν το πλήθος τους είναι \textbf{περιττό}, η διάμεσος είναι η \emph{μεσαία} τιμή. Αν το πλήθος τους είναι \textbf{άρτιο}, η διάμεσος είναι ο \emph{μέσος όρος των δύο μεσαίων} τιμών.
\end{kanonas}

Η διάμεσος χωρίζει τα δεδομένα στη μέση: τα μισά είναι μικρότερα ή ίσα με αυτήν και τα μισά μεγαλύτερα ή ίσα. Αν ένα δεδομένο γίνει πολύ μεγάλο, η διάμεσος σχεδόν δεν αλλάζει, ενώ η μέση τιμή «τραβιέται» προς αυτό.

\begin{kanonas}{μέση τιμή από πίνακα συχνοτήτων}
Όταν οι τιμές επαναλαμβάνονται, δεν προσθέτουμε μία μία. Εκτελούμε τρία βήματα:\\
1. Πολλαπλασιάζουμε κάθε τιμή με τη συχνότητά της.\\
2. Προσθέτουμε όλα τα γινόμενα.\\
3. Διαιρούμε με το σύνολο των παρατηρήσεων (το άθροισμα των συχνοτήτων).
\end{kanonas}

\begin{kanonas}{μέση τιμή από ομαδοποιημένα δεδομένα}
Όταν τα δεδομένα είναι χωρισμένα σε κλάσεις, δεν ξέρουμε την ακριβή τιμή κάθε παρατήρησης. Δεχόμαστε ότι όλες οι τιμές μιας κλάσης είναι ίσες με το \textbf{κέντρο της κλάσης} (το μέσο των άκρων της). Μετά κάνουμε τα ίδια τρία βήματα: κέντρο $\times$ συχνότητα, άθροισμα γινομένων, διαίρεση με το σύνολο. Το αποτέλεσμα είναι προσεγγιστικό.
\end{kanonas}

\begin{prosoxi}
\begin{itemize}[leftmargin=1.2em,itemsep=0.3ex,topsep=2pt]
\item Για τη διάμεσο \emph{πρώτα} βάζεις τις τιμές στη σειρά. Αν πάρεις την τιμή που κάθεται στη μέση της αταξινόμητης λίστας, κάνεις λάθος.
\item Στον πίνακα συχνοτήτων διαιρείς με το άθροισμα των \emph{συχνοτήτων}, όχι με το πλήθος των διαφορετικών τιμών.
\item Στα ομαδοποιημένα δεδομένα πολλαπλασιάζεις το \emph{κέντρο} της κλάσης με τη συχνότητα, όχι το αριστερό ή το δεξί άκρο.
\item Η διάμεσος δεν είναι ο μέσος όρος όλων των τιμών. Μόνο όταν το πλήθος είναι άρτιο παίρνουμε μέσο όρο, και μόνο των δύο μεσαίων.
\end{itemize}
\end{prosoxi}

\tmima{Λυμένα παραδείγματα}

\paradeigma{Η μέση τιμή}
\tixero Ο Νίκος σε πέντε αγώνες μπάσκετ πέτυχε $12$, $8$, $15$, $9$ και $21$ πόντους.\par
\tithelo Τη μέση τιμή των πόντων του.\par
\possindeo Προσθέτω τις πέντε τιμές και διαιρώ με το πλήθος τους, $5$.\par
\apantisi $\dfrac{12+8+15+9+21}{5}=\dfrac{65}{5}=13$ πόντοι ανά αγώνα.\par
\elegxos Το $13$ είναι ανάμεσα στο $8$ και στο $21$, όπως πρέπει. Και $5\cdot13=65$, που είναι το άθροισμα.

\paradeigma{Η διάμεσος, περιττό και άρτιο πλήθος}
\tixero Οι αναμονές της Ρήνης: $5, 6, 4, 5, 7, 5, 33$ λεπτά.\par
\tithelo Τη διάμεσο των επτά τιμών και τη διάμεσο αν χανόταν η καταγραφή της Δευτέρας ($5$), δηλαδή για τις $6, 4, 5, 7, 5, 33$.\par
\possindeo Βάζω τις τιμές στη σειρά. Επτά τιμές: μεσαία είναι η τέταρτη. Έξι τιμές: μέσος όρος της τρίτης και της τέταρτης.\par
\apantisi Επτά τιμές: $4, 5, 5, \underline{5}, 6, 7, 33$, άρα διάμεσος $5$ λεπτά. Έξι τιμές: $4, 5, \underline{5}, \underline{6}, 7, 33$, άρα διάμεσος $\dfrac{5+6}{2}=5{,}5$ λεπτά.\par
\elegxos Η μέση τιμή των επτά είναι $\dfrac{65}{7}\approx 9{,}3$ και την τραβά προς τα πάνω η τιμή $33$. Η διάμεσος $5$ δείχνει καλύτερα την «συνηθισμένη» αναμονή.

\paradeigma{Ποια τιμή λείπει;}
\tixero Ο Ηλίας έγραψε στα Μαθηματικά $15$, $17$, $14$, $18$ σε τέσσερα τεστ και θέλει μέσο όρο $16$ μετά το πέμπτο.\par
\tithelo Τι βαθμό χρειάζεται στο πέμπτο τεστ;\par
\possindeo Αν η μέση τιμή των $5$ τεστ είναι $16$, το άθροισμά τους είναι $5\cdot16=80$.\par
\apantisi Τα τέσσερα τεστ δίνουν $15+17+14+18=64$, άρα στο πέμπτο χρειάζεται $80-64=16$.\par
\elegxos $\dfrac{15+17+14+18+16}{5}=\dfrac{80}{5}=16$.

\paradeigma{Πίνακας συχνοτήτων}
\tixero Η Φανή ρώτησε $25$ συμμαθητές πόσα αδέλφια έχουν.\par
\tithelo Τη μέση τιμή.\par
\possindeo Προσθέτω μια στήλη «τιμή $\times$ συχνότητα».\par

\begin{center}
\renewcommand{\arraystretch}{1.25}
\begin{tabular}{lccccc}
\toprule
Αδέλφια & $0$ & $1$ & $2$ & $3$ & $4$\\
Συχνότητα & $3$ & $9$ & $8$ & $4$ & $1$\\
Τιμή $\times$ συχνότητα & $0$ & $9$ & $16$ & $12$ & $4$\\
\bottomrule
\end{tabular}
\end{center}

\apantisi Σύνολο παρατηρήσεων $3+9+8+4+1=25$. Άθροισμα γινομένων $0+9+16+12+4=41$. Μέση τιμή $\dfrac{41}{25}=1{,}64$ αδέλφια.\par
\elegxos Το $1{,}64$ είναι ανάμεσα στο $0$ και στο $4$. Δεν υπάρχει μαθητής με $1{,}64$ αδέλφια, αλλά ως μέσος όρος έχει νόημα.

\paradeigma{Ομαδοποιημένα δεδομένα}
\tixero Ο χρόνος που κάνουν $25$ μαθητές να φτάσουν στο σχολείο ομαδοποιήθηκε σε κλάσεις.\par
\tithelo Τον μέσο χρόνο διαδρομής.\par
\possindeo Βρίσκω το κέντρο κάθε κλάσης (π.χ. για $0$--$5$ είναι $\frac{0+5}{2}=2{,}5$) και συνεχίζω όπως στον πίνακα συχνοτήτων.\par

\begin{center}
\renewcommand{\arraystretch}{1.25}
\begin{tabular}{lccccc}
\toprule
Χρόνος (min) & $0$--$5$ & $5$--$10$ & $10$--$15$ & $15$--$20$ & $20$--$25$\\
Κέντρο κλάσης & $2{,}5$ & $7{,}5$ & $12{,}5$ & $17{,}5$ & $22{,}5$\\
Συχνότητα & $4$ & $9$ & $7$ & $3$ & $2$\\
Κέντρο $\times$ συχνότητα & $10$ & $67{,}5$ & $87{,}5$ & $52{,}5$ & $45$\\
\bottomrule
\end{tabular}
\end{center}

\apantisi Σύνολο $4+9+7+3+2=25$ μαθητές. Άθροισμα γινομένων $10+67{,}5+87{,}5+52{,}5+45=262{,}5$. Μέση τιμή $\dfrac{262{,}5}{25}=10{,}5$ λεπτά.\par
\elegxos Η μέση τιμή είναι ανάμεσα στο $0$ και στο $25$ και πλησιάζει τις κλάσεις με τις περισσότερες παρατηρήσεις ($5$--$10$ και $10$--$15$). Είναι προσέγγιση, γιατί δεν ξέρουμε τις ακριβείς τιμές.

\tmima{Ασκήσεις}

\zoni{Το κατάλαβα}{μέση τιμή, διάμεσος, πίνακας}

\ex Βρες τη μέση τιμή: α) $24$ παρατηρήσεις με άθροισμα $120$. β) $8$ παρατηρήσεις με άθροισμα $20$. γ) $40$ παρατηρήσεις με άθροισμα $10$.

\ex Βρες το άθροισμα: α) $30$ παρατηρήσεων με μέση τιμή $4$. β) $6$ παρατηρήσεων με μέση τιμή $2{,}5$. γ) $20$ παρατηρήσεων με μέση τιμή $0{,}1$.

\ex Βρες το πλήθος των παρατηρήσεων που έχουν: α) μέση τιμή $4$ και άθροισμα $48$. β) μέση τιμή $2{,}5$ και άθροισμα $30$. γ) μέση τιμή $0{,}5$ και άθροισμα $20$.

\ex Υπολόγισε τη μέση τιμή: α) $4, 6, 8, 10, 12$. β) $-2, 0, 3, 5$ (θερμοκρασίες σε $^\circ$C). γ) $7, 7, 7, 7$.

\ex Βρες τη διάμεσο: α) $9, 3, 7, 1, 5$. β) $12, 4, 8, 6$. γ) $10, 3, 3, 7, 3, 9$.

\ex Σωστό ή λάθος; α) Η μέση τιμή είναι πάντα μία από τις τιμές των δεδομένων. β) Αν τα δεδομένα είναι έξι, η διάμεσος είναι ο μέσος όρος της τρίτης και της τέταρτης τιμής, μετά τη σειρά. γ) Για να βρεις τη διάμεσο, πρώτα βάζεις τα δεδομένα σε αύξουσα σειρά. δ) Αν όλες οι τιμές είναι ακέραιοι, η μέση τιμή είναι κι αυτή ακέραιος. ε) Η μέση τιμή μπορεί να είναι μεγαλύτερη από τη μεγαλύτερη τιμή.

\ex Συμπλήρωσε τον πίνακα και βρες τη μέση τιμή: \par
\begin{center}
\renewcommand{\arraystretch}{1.25}
\begin{tabular}{lcccc}
\toprule
Τιμή & $0$ & $1$ & $2$ & $3$\\
Συχνότητα & $4$ & $6$ & $8$ & $2$\\
Τιμή $\times$ συχνότητα & $\ldots$ & $\ldots$ & $\ldots$ & $\ldots$\\
\bottomrule
\end{tabular}
\end{center}

\zoni{Το χρησιμοποιώ}{σπορ, χωριό, σχολείο}

\ex \textbf{Οι βολές.} Ο Νίκος σε έξι προπονήσεις έβαλε $6$, $7$, $5$, $8$, $9$ και $7$ βολές από τις $10$. Βρες τη μέση τιμή και τη διάμεσο.

\ex \textbf{Τα τριαντάφυλλα.} Η θεία Μαρία έχει $9$ τριανταφυλλιές με μέσο όρο $12$ τριαντάφυλλα η καθεμία. α) Πόσα τριαντάφυλλα έχουν όλες μαζί; β) Μια δέκατη τριανταφυλλιά έχει $22$ τριαντάφυλλα. Ποιος είναι ο νέος μέσος όρος;

\ex \textbf{Ο βαθμός που λείπει.} Η Ζωή έγραψε $16$, $18$, $15$, $17$ στα τέσσερα πρώτα τεστ. Τι βαθμό χρειάζεται στο πέμπτο ώστε ο μέσος όρος των πέντε να είναι $17$;

\ex \textbf{Τα βιβλία του καλοκαιριού.} Η Αγγελική ρώτησε $24$ παιδιά πόσα βιβλία διάβασαν το καλοκαίρι.\par
\begin{center}
\renewcommand{\arraystretch}{1.25}
\begin{tabular}{lcccccc}
\toprule
Βιβλία & $0$ & $1$ & $2$ & $3$ & $4$ & $5$\\
Συχνότητα & $1$ & $4$ & $7$ & $7$ & $4$ & $1$\\
\bottomrule
\end{tabular}
\end{center}
Βρες τη μέση τιμή.

\ex \textbf{Το χαρτζιλίκι.} Ο Στέλιος ρώτησε $20$ συμμαθητές πόσα ευρώ παίρνουν την εβδομάδα.\par
\begin{center}
\renewcommand{\arraystretch}{1.25}
\begin{tabular}{lcccc}
\toprule
Ευρώ & $0$--$5$ & $5$--$10$ & $10$--$15$ & $15$--$20$\\
Συχνότητα & $3$ & $8$ & $6$ & $3$\\
\bottomrule
\end{tabular}
\end{center}
Βρες τα κέντρα των κλάσεων και τη μέση τιμή.

\zoni{Το σκέφτομαι}{συγκρίσεις και αιτιολογήσεις}

\ex Πέντε δεδομένα έχουν μέση τιμή $10$. Προσθέτουμε ένα έκτο δεδομένο ίσο με $10$. Πώς αλλάζει η μέση τιμή; Αν το έκτο ήταν $16$, ποια θα ήταν;

\ex Βρες πέντε αριθμούς που να έχουν διάμεσο $7$ και μέση τιμή $10$. Είναι μοναδική η απάντηση;

\ex Το τμήμα Α έχει $10$ μαθητές με μέσο όρο $14$ και το τμήμα Β έχει $20$ μαθητές με μέσο όρο $17$. Ο Στέλιος λέει: «Ο μέσος όρος και των δύο μαζί είναι $\frac{14+17}{2}=15{,}5$». Έχει δίκιο; Βρες τον σωστό μέσο όρο.

\ex Σε ένα μαγαζί του χωριού πέντε εργαζόμενοι παίρνουν $900$, $1000$, $1100$, $1000$ και $5000$ ευρώ τον μήνα (ο τελευταίος είναι ο ιδιοκτήτης). Βρες τη μέση τιμή και τη διάμεσο. Ποια από τις δύο περιγράφει καλύτερα τον «τυπικό» μισθό και γιατί;

\zoni{Για το σπίτι}{μετρήσεις στο σπίτι}

\ex Κράτησε σημειώσεις για τις ώρες που κοιμήθηκες επτά βραδιές. Βρες τη μέση τιμή και τη διάμεσο.

\ex Ρώτησε $10$ ανθρώπους πόσα κινητά τηλέφωνα έχουν στο σπίτι τους. Φτιάξε πίνακα συχνοτήτων και βρες τη μέση τιμή.

\paragrafos{Εύρος, κορυφές και απόμακρες τιμές}{Γ.2.2, «Εύρος, κορυφές, πολλαπλές κορυφές, απόμακρες τιμές / Μεταβλητότητα δεδομένων»}

\begin{anagke}
Ο Νίκος και η Ρήνη παίζουν στο διάλειμμα ένα παιχνίδι με βολές στο καλάθι. Κάνουν πέντε γύρους των $10$ βολών ο καθένας. Ο Νίκος έβαλε $1$, $9$, $6$, $10$ και $4$ καλάθια. Η Ρήνη έβαλε $5$, $6$, $6$, $6$ και $7$.

«Ίδιος μέσος όρος!», φωνάζει ο Νίκος, «και οι δύο έχουμε $6$». Η Ρήνη κουνά το κεφάλι: «Εγώ ξέρω σχεδόν πάντα τι θα κάνω. Εσύ μια παίρνεις δέκα και μια ένα. Δεν παίζουμε το ίδιο». Η κυρία Αλεξάνδρα ρωτά: «Πώς θα έδειχνες με έναν αριθμό αυτό που βλέπει η Ρήνη;».

Η μέση τιμή δεν φτάνει για να περιγράψουμε δεδομένα. Χρειαζόμαστε και ένα μέτρο για το πόσο «απλώνονται» οι τιμές.

Σε αυτή την παράγραφο μαθαίνουμε το εύρος, την κορυφή (και τις πολλαπλές κορυφές), τις απόμακρες τιμές και τι σημαίνει μεταβλητότητα των δεδομένων.
\end{anagke}

\begin{lexi}{εύρος, κορυφή, απόμακρη τιμή, μεταβλητότητα}{ονόματα και σημασίες}
Το \lexiword{εύρος} είναι η έκταση (το «πλάτος») που καταλαμβάνουν τα δεδομένα. Η \lexiword{κορυφή} είναι η κορυφή ενός λόφου: η τιμή που «ξεπροβάλλει» στο διάγραμμα επειδή εμφανίζεται πιο συχνά (και λέγεται και \emph{επικρατούσα τιμή}). Η \lexiword{απόμακρη} τιμή βρίσκεται μακριά από τις άλλες. Η \lexiword{μεταβλητότητα} είναι το πόσο μεταβάλλονται, δηλαδή πόσο διαφέρουν μεταξύ τους, οι τιμές.
\end{lexi}

\begin{metafrasi}
Η μεγαλύτερη τιμή μείον τη μικρότερη. & $R=\text{μεγαλύτερη}-\text{μικρότερη}$\\
Η τιμή που εμφανίζεται τις περισσότερες φορές. & κορυφή (επικρατούσα τιμή)\\
Δύο ή περισσότερες τιμές με την ίδια, μέγιστη, συχνότητα. & πολλαπλές κορυφές\\
Μια τιμή που ξεχωρίζει πολύ από τις υπόλοιπες. & απόμακρη τιμή\\
Οι τιμές είναι πολύ διαφορετικές μεταξύ τους. & μεγάλη μεταβλητότητα\\
\end{metafrasi}

\tmima{Η θεωρία}

\begin{kanonas}{εύρος}
$$\text{εύρος }R=(\text{μεγαλύτερη τιμή})-(\text{μικρότερη τιμή})$$
Σε \textbf{ομαδοποιημένα} δεδομένα το εύρος είναι η διαφορά του μικρότερου άκρου της πρώτης κλάσης από το μεγαλύτερο άκρο της τελευταίας κλάσης.
\end{kanonas}

Στο παιχνίδι του Νίκου και της Ρήνης, τα δεδομένα στην αριθμογραμμή δείχνουν καθαρά τη διαφορά.

\begin{center}
\begin{tikzpicture}[x=0.8cm]
\draw[-{Stealth},thick] (-0.5,0)--(11,0);
\foreach \x in {0,...,10}{\draw (\x,0.1)--(\x,-0.1) node[below,font=\scriptsize]{$\x$};}
\foreach \x in {1,4,6,9,10}{\fill[gwarmD] (\x,0.5) circle (3.5pt);}
\foreach \x in {5,7}{\fill[gblue] (\x,1.2) circle (3.5pt);}
\fill[gblue] (6,1.2) circle (3.5pt);\fill[gblue] (6,1.55) circle (3.5pt);\fill[gblue] (6,1.9) circle (3.5pt);
\node[font=\small,gwarmD,anchor=west] at (11.2,0.5) {Νίκος: $R=9$};
\node[font=\small,gblue,anchor=west] at (11.2,1.4) {Ρήνη: $R=2$};
\end{tikzpicture}
\end{center}

\begin{kanonas}{κορυφή}
\textbf{Κορυφή} (ή \textbf{επικρατούσα τιμή}) είναι η τιμή που εμφανίζεται τις περισσότερες φορές, δηλαδή έχει τη μεγαλύτερη συχνότητα. Αν δύο ή περισσότερες τιμές έχουν την ίδια μεγαλύτερη συχνότητα, έχουμε \textbf{πολλαπλές κορυφές}.
\end{kanonas}

Η κορυφή έχει νόημα και σε δεδομένα που δεν είναι αριθμοί. Αν ρωτήσουμε ποιο είναι το αγαπημένο χρώμα, η κορυφή μπορεί να είναι «μπλε», ενώ μέση τιμή χρωμάτων δεν υπάρχει. Σε ιστόγραμμα κορυφή είναι η κλάση με το ψηλότερο ορθογώνιο. Το ιστόγραμμα εδώ έχει δύο κορυφές:

\begin{center}
\binhist{2,7,3,6,1}{0,10,20,30,40,50}{8}{διάρκεια διαδρομής (λεπτά)}
\end{center}

\begin{orismos}{απόμακρη τιμή}
\textbf{Απόμακρη τιμή} λέγεται το δεδομένο που διαφέρει σημαντικά από τα υπόλοιπα δεδομένα.
\end{orismos}

Πού οφείλονται οι απόμακρες τιμές; Συχνά σε \emph{λάθος}: κάποιος πληκτρολόγησε $85$ αντί για $8{,}5$ ή το όργανο μέτρησης δεν δούλευε σωστά. Άλλες φορές είναι \emph{φυσική ποικιλία}: ανάμεσα σε παιδιά σε μια παιδική χαρά, ο παππούς είναι «απόμακρος» στην ηλικία, χωρίς να υπάρχει λάθος. Δεν υπάρχει αριθμητικός τύπος που να μας λέει πότε μια τιμή είναι απόμακρη: κοιτάμε τα δεδομένα και ρωτάμε αν μια τιμή «ξεκόβει» από τις άλλες.

Η απόμακρη τιμή μεγαλώνει πολύ το εύρος και τραβά τη μέση τιμή, ενώ η διάμεσος και η κορυφή αλλάζουν λίγο ή καθόλου.

\begin{idiotita}{μεταβλητότητα}
Όσο μεγαλύτερο είναι το εύρος, τόσο περισσότερο απλώνονται τα δεδομένα και τόσο μεγαλύτερη είναι η \textbf{μεταβλητότητά} τους. Για να συγκρίνουμε δύο σύνολα δεδομένων κοιτάμε μαζί τη μέση τιμή (πού βρίσκονται) και το εύρος (πόσο απλώνονται). Δύο σύνολα μπορεί να έχουν την ίδια μέση τιμή και πολύ διαφορετικό εύρος.
\end{idiotita}

\begin{prosoxi}
\begin{itemize}[leftmargin=1.2em,itemsep=0.3ex,topsep=2pt]
\item Η κορυφή είναι η \emph{τιμή}, όχι η συχνότητά της. Στα δεδομένα $2,5,5,5,5,7$ η κορυφή είναι το $5$ και όχι το $4$ (το $4$ είναι πόσες φορές εμφανίζεται).
\item Το εύρος είναι αριθμός που δεν είναι ποτέ αρνητικός. Είναι πάντα μεγαλύτερο μείον μικρότερο, και ας είναι οι τιμές αρνητικές: για $-3$ και $4$ βγαίνει $4-(-3)=7$.
\item Δεν πετάμε μια απόμακρη τιμή χωρίς λόγο. Πρώτα ρωτάμε από πού προέρχεται, αν είναι λάθος ή αληθινή.
\item Το εύρος κοιτά μόνο τα δύο άκρα. Μια απόμακρη τιμή αρκεί για να το μεγαλώσει πολύ.
\end{itemize}
\end{prosoxi}

\tmima{Λυμένα παραδείγματα}

\paradeigma{Το εύρος}
\tixero Η Πολυξένη κατέγραψε στη Θεσσαλονίκη τις θερμοκρασίες επτά ημερών ($^\circ$C): $-3$, $0$, $2$, $-1$, $4$, $1$, $2$.\par
\tithelo Το εύρος.\par
\possindeo Βρίσκω τη μεγαλύτερη και τη μικρότερη τιμή.\par
\apantisi Μεγαλύτερη $4$, μικρότερη $-3$. Εύρος $R=4-(-3)=7$ $^\circ$C.\par
\elegxos Από το $-3$ ως το $4$ στην αριθμογραμμή υπάρχουν $7$ μονάδες.

\paradeigma{Η κορυφή και οι πολλαπλές κορυφές}
\tixero Τα νούμερα παπουτσιών $12$ παιδιών: $36$, $37$, $37$, $38$, $36$, $39$, $37$, $38$, $36$, $40$, $38$, $37$.\par
\tithelo Την κορυφή.\par
\possindeo Μετρώ πόσες φορές εμφανίζεται κάθε τιμή.\par
\apantisi Το $36$ εμφανίζεται $3$ φορές, το $37$ εμφανίζεται $4$ φορές, το $38$ εμφανίζεται $3$ φορές, τα $39$ και $40$ από μία. Κορυφή το $37$.\par
\elegxos Οι συχνότητες αθροίζουν $3+4+3+1+1=12$, όσα και τα παιδιά. Στα δεδομένα $5,6,6,7,8,8,9$ τα $6$ και $8$ εμφανίζονται από δύο φορές και τα άλλα από μία: υπάρχουν δύο κορυφές, το $6$ και το $8$.

\paradeigma{Απόμακρη τιμή}
\tixero Σε ένα οικογενειακό τραπέζι οι ηλικίες των παιδιών είναι $9, 10, 11, 11, 12, 12, 12$ και ο παππούς Γιάννης είναι $83$.\par
\tithelo Να δούμε πώς επηρεάζει ο παππούς τη μέση τιμή, τη διάμεσο και το εύρος.\par
\possindeo Υπολογίζω πρώτα χωρίς τον παππού και μετά με τον παππού.\par
\apantisi Χωρίς τον παππού: άθροισμα $77$, μέση τιμή $77:7=11$, διάμεσος $11$, εύρος $12-9=3$. Με τον παππού: άθροισμα $160$, μέση τιμή $160:8=20$, διάμεσος $\dfrac{11+12}{2}=11{,}5$, εύρος $83-9=74$.\par
\elegxos Η μέση τιμή σχεδόν διπλασιάστηκε και το εύρος έγινε $74$, ενώ η διάμεσος αλλάζει ελάχιστα. Η μέση ηλικία $20$ δεν περιγράφει κανέναν στο τραπέζι. Η $83$ είναι απόμακρη, αλλά όχι λάθος: είναι ο παππούς.

\paradeigma{Ίδιος μέσος όρος, διαφορετική μεταβλητότητα}
\tixero Βολές του Νίκου: $1,9,6,10,4$. Βολές της Ρήνης: $5,6,6,6,7$.\par
\tithelo Να συγκρίνουμε τους δύο παίκτες.\par
\possindeo Υπολογίζω μέση τιμή και εύρος για τον καθένα.\par
\apantisi Νίκος: μέση τιμή $\frac{30}{5}=6$, εύρος $10-1=9$. Ρήνη: μέση τιμή $\frac{30}{5}=6$, εύρος $7-5=2$. Και οι δύο έχουν $6$ κατά μέσο όρο, αλλά η Ρήνη είναι πολύ πιο σταθερή (μικρό εύρος), ενώ ο Νίκος είναι απρόβλεπτος (μεγάλο εύρος).\par
\elegxos Το σύνολο της Ρήνης είναι «μαζεμένο» γύρω από το $6$ και του Νίκου απλωμένο από το $1$ ως το $10$, όπως φαίνεται στην αριθμογραμμή.

\paradeigma{Εύρος σε ομαδοποιημένα δεδομένα}
\tixero Τα ύψη $20$ παιδιών (cm): $140$--$145$: $2$ παιδιά, $145$--$150$: $5$, $150$--$155$: $8$, $155$--$160$: $4$, $160$--$165$: $1$.\par
\tithelo Το εύρος και την κλάση με τη μεγαλύτερη συχνότητα.\par
\possindeo Εύρος: μεγαλύτερο άκρο της τελευταίας κλάσης μείον μικρότερο άκρο της πρώτης.\par
\apantisi Εύρος $165-140=25$ cm. Το ψηλότερο ορθογώνιο είναι της κλάσης $150$--$155$ ($8$ παιδιά).\par
\elegxos Το εύρος είναι προσεγγιστικό, γιατί δεν ξέρουμε το πραγματικό ύψος του ψηλότερου ή του πιο κοντού παιδιού.

\tmima{Ασκήσεις}

\zoni{Το κατάλαβα}{εύρος, κορυφή, απόμακρη τιμή}

\ex Βρες το εύρος: α) $3,8,5,12,7$. β) $-4,0,6,-2$. γ) $15,15,15$. δ) $0{,}5$, $2{,}5$, $1{,}75$.

\ex Βρες την κορυφή ή τις κορυφές: α) $2,3,3,5,3,7$. β) $1,4,4,6,6,9$. γ) κόκκινο, μπλε, μπλε, πράσινο, κόκκινο, μπλε.

\ex Σωστό ή λάθος; α) Το εύρος μπορεί να είναι αρνητικός αριθμός. β) Αν το εύρος είναι $0$, όλες οι τιμές είναι ίσες. γ) Η κορυφή είναι πάντα μία. δ) Η κορυφή είναι η τιμή που εμφανίζεται τις περισσότερες φορές. ε) Μια απόμακρη τιμή μεγαλώνει το εύρος.

\ex Βρες, αν υπάρχει, την απόμακρη τιμή: α) $4,5,6,5,4,48$. β) $120,118,125,3,122$. γ) $10,12,11,13,12$.

\ex Το ιστόγραμμα δείχνει πόσοι μαθητές της Α3 μένουν σε κάθε απόσταση από το σχολείο (σε λεπτά με τα πόδια).
\begin{center}
\binhist{2,7,3,6,1}{0,10,20,30,40,50}{8}{χρόνος (λεπτά)}
\end{center}
α) Πόσοι μαθητές ρωτήθηκαν; β) Ποιο είναι το εύρος; γ) Έχει το ιστόγραμμα περισσότερες από μία κορυφές; δ) Πόσοι μαθητές χρειάζονται λιγότερο από $20$ λεπτά;

\ex Οι θερμοκρασίες μιας εβδομάδας στο χωριό ($^\circ$C), από Δευτέρα ως Κυριακή: $-2, 3, 0, 5, 2, 8, -2$. Βρες τη μέση τιμή, το εύρος και την κορυφή.

\ex Δύο σύνολα δεδομένων: Α: $5,6,7,8,9$ και Β: $1,4,7,10,13$. Βρες τη μέση τιμή και το εύρος του καθενός. Ποιο έχει μεγαλύτερη μεταβλητότητα;

\zoni{Το χρησιμοποιώ}{διαγωνισμοί, πωλήσεις, διαφημίσεις}

\ex \textbf{Οι δρομείς.} Σε πέντε αγώνες των $100$ m η Αγγελική έκανε $16, 17, 17, 17, 18$ δευτερόλεπτα και η Πολυξένη $14, 15, 17, 19, 20$. Βρες τη μέση τιμή και το εύρος της καθεμιάς. Ποια είναι πιο σταθερή; Ποια έκανε τον καλύτερο χρόνο;

\ex \textbf{Τα παγωτά.} Σε ένα περίπτερο στην παραλία πουλήθηκαν επτά μέρες $40, 36, 44, 38, 42, 40$ και $5$ παγωτά. Την τελευταία μέρα έβρεχε. α) Βρες μέση τιμή, διάμεσο και εύρος και με τις επτά μέρες και χωρίς τη βροχερή. β) Ποιο μέτρο άλλαξε λιγότερο;

\ex \textbf{Τα ύψη.} Στον πίνακα φαίνονται τα ύψη $20$ παιδιών:\par
\begin{center}
\renewcommand{\arraystretch}{1.25}
\begin{tabular}{lccccc}
\toprule
Ύψος (cm) & $140$--$145$ & $145$--$150$ & $150$--$155$ & $155$--$160$ & $160$--$165$\\
Παιδιά & $2$ & $5$ & $8$ & $4$ & $1$\\
\bottomrule
\end{tabular}
\end{center}
Βρες το εύρος και τη μέση τιμή του ύψους (με τα κέντρα των κλάσεων).

\ex \textbf{Οι βαθμοί.} Ο πίνακας δείχνει τους βαθμούς $22$ μαθητών σε ένα τεστ (κλίμακα $0$--$20$):\par
\begin{center}
\renewcommand{\arraystretch}{1.25}
\begin{tabular}{lcccccc}
\toprule
Βαθμός & $12$ & $13$ & $14$ & $15$ & $16$ & $17$\\
Μαθητές & $2$ & $3$ & $6$ & $6$ & $3$ & $2$\\
\bottomrule
\end{tabular}
\end{center}
Βρες την κορυφή (ή τις κορυφές), το εύρος και τη μέση τιμή.

\ex \textbf{Το ζάρι.} Ο Ηλίας έριξε ένα ζάρι $12$ φορές και έγραψε: $3,5,2,6,5,1,5,4,6,5,2,3$. Φτιάξε πίνακα συχνοτήτων και βρες την κορυφή, το εύρος και τη μέση τιμή.

\zoni{Το σκέφτομαι}{τι αλλάζει και τι μένει}

\ex Βρες πέντε αριθμούς με εύρος $10$, διάμεσο $7$ και κορυφή $7$.

\ex Τα δεδομένα $2, 5, 8$ έχουν μέση τιμή $5$ και εύρος $6$. α) Πρόσθεσε $3$ σε κάθε τιμή. Τι γίνεται η μέση τιμή και τι το εύρος; β) Πολλαπλασίασε κάθε τιμή με $2$. Τι γίνεται τώρα; Τι παρατηρείς;

\ex Τα δεδομένα $3,4,5,5,6,7$ έχουν μέση τιμή $5$ και διάμεσο $5$. Πρόσθεσε το $40$. Βρες τη νέα μέση τιμή και τη νέα διάμεσο. Τι κάνει η απόμακρη τιμή σε καθένα από τα δύο;

\ex Στο σχολικό πανηγύρι ρωτήσαμε τις ηλικίες όσων ήρθαν.
\begin{center}
\binhist{2,7,1,3,6}{0,10,20,30,40,50}{8}{ηλικία (έτη)}
\end{center}
α) Πόσες κορυφές έχει το ιστόγραμμα; β) Πώς το εξηγείς; γ) Βρες τη μέση τιμή (με τα κέντρα των κλάσεων). Περιγράφει κάποιον από τους παρευρισκόμενους;

\zoni{Για το σπίτι}{μέτρησε και σύγκρινε}

\ex Μέτρησε επί πέντε μέρες πόσα λεπτά κάνεις να φτάσεις στο σχολείο. Βρες τη μέση τιμή, τη διάμεσο και το εύρος.

\ex Ρώτησε $12$ ανθρώπους ποιο είναι το αγαπημένο τους χρώμα. Βρες την κορυφή. Γιατί δεν έχει νόημα να βρεις μέση τιμή;

\epanalipsi{Γ2}

\tmima{Λεξικό του κεφαλαίου}

\renewcommand{\arraystretch}{1.5}
\begin{longtable}{@{}>{\raggedright\arraybackslash}p{0.28\linewidth} >{\raggedright\arraybackslash}p{0.62\linewidth}@{}}
\toprule
\textbf{Λέξη ή σύμβολο} & \textbf{Με μία πρόταση}\\
\midrule
\endhead
μέση τιμή & Άθροισμα των δεδομένων διά το πλήθος τους.\\
διάμεσος & Η μεσαία τιμή στη σειρά (ή ο μέσος όρος των δύο μεσαίων).\\
μέση τιμή από πίνακα & Άθροισμα των γινομένων τιμή $\times$ συχνότητα, διά το σύνολο των παρατηρήσεων.\\
κέντρο κλάσης & Το μέσο των δύο άκρων της κλάσης.\\
εύρος $R$ & Μεγαλύτερη τιμή μείον μικρότερη τιμή.\\
κορυφή & Η τιμή με τη μεγαλύτερη συχνότητα (μπορεί να υπάρχουν πολλές).\\
απόμακρη τιμή & Τιμή που διαφέρει σημαντικά από τις υπόλοιπες.\\
μεταβλητότητα & Πόσο απλώνονται και διαφέρουν οι τιμές μεταξύ τους.\\
\bottomrule
\end{longtable}

\tmima{Ασκήσεις επανάληψης}

\ex Βρες τη μέση τιμή των $7, 9, 12, 4, 8$.

\ex Βρες τη διάμεσο των $11, 3, 8, 6, 15$ και των $7, 2, 9, 4$.

\ex Δώδεκα παρατηρήσεις έχουν μέση τιμή $6{,}5$. Πόσο είναι το άθροισμά τους;

\ex Η Ζωή έγραψε $13, 16, 14, 15$ στα τέσσερα πρώτα διαγωνίσματα. Τι βαθμό χρειάζεται στο πέμπτο ώστε ο μέσος όρος των πέντε να είναι $15$;

\ex Ο πίνακας δείχνει πόσες ώρες γυμναστικής κάνουν την εβδομάδα $20$ μαθητές. Βρες τη μέση τιμή.\par
\begin{center}
\renewcommand{\arraystretch}{1.25}
\begin{tabular}{lcccc}
\toprule
Ώρες & $1$ & $2$ & $3$ & $4$\\
Μαθητές & $5$ & $8$ & $4$ & $3$\\
\bottomrule
\end{tabular}
\end{center}

\ex Στον πίνακα φαίνονται τα λεπτά που περνούν $20$ παιδιά στο κινητό κάθε μέρα. Βρες τη μέση τιμή με τα κέντρα των κλάσεων.\par
\begin{center}
\renewcommand{\arraystretch}{1.25}
\begin{tabular}{lccc}
\toprule
Λεπτά & $0$--$20$ & $20$--$40$ & $40$--$60$\\
Παιδιά & $4$ & $10$ & $6$\\
\bottomrule
\end{tabular}
\end{center}

\ex Οι ηλικίες επτά αθλητών είναι $14, 18, 14, 22, 17, 14, 19$. Βρες το εύρος και την κορυφή.

\ex Οι βαθμοί έξι μαθητών είναι $15, 16, 14, 15, 2, 16$. Ποια είναι η απόμακρη τιμή; Βρες τη μέση τιμή με και χωρίς αυτή.

\ex Δύο παίκτες πέτυχαν πόντους σε πέντε αγώνες. Ο πρώτος: $2, 10, 6, 14, 8$. Ο δεύτερος: $7, 8, 9, 8, 8$. Βρες τη μέση τιμή και το εύρος του καθενός και σχολίασε.

\ex Το ιστόγραμμα δείχνει πόσες ώρες την εβδομάδα διαβάζουν $20$ μαθητές.
\begin{center}
\binhist{1,4,7,5,3}{0,2,4,6,8,10}{8}{ώρες την εβδομάδα}
\end{center}
Βρες το εύρος, την κλάση-κορυφή και τη μέση τιμή.

\ex Δέκα αριθμοί έχουν μέση τιμή $8$. Προσθέτουμε τον $19$. Ποια είναι η νέα μέση τιμή;

\ex Στο τμήμα Α φοιτούν $5$ μαθητές με μέσο όρο $12$ και στο τμήμα Β $15$ μαθητές με μέσο όρο $16$. Ποιος είναι ο μέσος όρος και των δύο μαζί;

\ex Βρες έξι αριθμούς με διάμεσο $5{,}5$, εύρος $8$ και μέση τιμή $6$.

\ex Οι θερμοκρασίες έξι ημερών ($^\circ$C) είναι $-1, 2, -3, 4, 0, 2$. Βρες μέση τιμή, διάμεσο, εύρος και κορυφή.

\ex Σωστό ή λάθος; α) Η διάμεσος επτά αριθμών είναι ο μέσος όρος των τριών μεσαίων. β) Τέσσερις αριθμοί με άθροισμα $20$ έχουν μέση τιμή $5$. γ) Τα δεδομένα $3,3,5,5,7$ έχουν δύο κορυφές. δ) Οι απόμακρες τιμές δεν επηρεάζουν τη μέση τιμή. ε) Το εύρος των $-5, 0, 4$ είναι $1$.

\ex Σωστό ή λάθος; α) Δύο σύνολα με την ίδια μέση τιμή έχουν πάντα το ίδιο εύρος. β) Όσο μεγαλύτερο το εύρος, τόσο μεγαλύτερη η μεταβλητότητα. γ) Σε ομαδοποιημένα δεδομένα το εύρος είναι το μεγαλύτερο άκρο της τελευταίας κλάσης μείον το μικρότερο άκρο της πρώτης. δ) Αν προσθέσουμε $5$ σε όλες τις τιμές, το εύρος αυξάνεται κατά $5$. ε) Η κορυφή μπορεί να είναι ένα χρώμα.

\tmima{Χάρτης επανάληψης}

\begin{center}
\renewcommand{\arraystretch}{1.3}
\begin{tabular}{ll}
\toprule
\textbf{Ασκήσεις} & \textbf{Παράγραφος}\\
\midrule
1, 2, 3, 4, 5, 6, 11, 12 & Γ2.1 Μέση τιμή και διάμεσος\\
7, 8, 9, 10, 13, 14 & Γ2.2 Εύρος, κορυφές και απόμακρες τιμές\\
15, 16 & Γ2.1 και Γ2.2\\
\bottomrule
\end{tabular}
\end{center}

